% ============================================================
%  Chapter 8 — The Cloud
%  The Secret Life of the Internet
% ============================================================

\chapter{The Cloud}

\chapteropening{%
  People say things are ``in the cloud''.
  But what does that actually mean?
  Cloud Pal is here to reveal the surprising truth!
}
\chapterbadge{Cloud Clues Unlocked}

% --- Opening story ---
\section*{The Reveal}

Sam had drawn a picture on the tablet.
The tablet said: ``Saved to the cloud.''

Sam looked out the window at the grey, rainy sky.

\begin{center}
  \Large\itshape\color{InternetBlue}
  ``Is my drawing floating up there somewhere?''
\end{center}

A fluffy, smiling cloud shape floated over.
But then it unzipped itself — and inside was a row of servers, blinking lights and all.

\character{Cloud Pal}{Surprise!  I'm not actually a sky cloud.
  That is just the nickname people gave me because I seem magical and hard to see.
  Really, I am just computers — lots and lots of computers — in buildings far away.
  And right now, your drawing is safely stored on one of them!}

\sectionrule

% --- What is the cloud? ---
\section*{What Is the Cloud?}

\term{The cloud} is a way of storing, accessing, and using files, programmes, and computing power
on servers located far away from your device — via the internet.

\begin{funfact}
Every minute, people upload around 500 hours of video to YouTube, share 695,000 Instagram Stories, and send 16 million text messages — all stored in the cloud!
\end{funfact}

Instead of keeping everything on your own tablet or laptop,
the cloud lets you keep things on servers somewhere else.
As long as you have an internet connection, you can get to your stuff
from any device, anywhere in the world.

\begin{picturethis}
  Imagine you kept all your favourite toys in a giant warehouse instead of your bedroom.
  Any time you wanted to play with one, you could ask the warehouse to send it to you.
  When you visited your grandparents and wanted your favourite toy,
  you could get it there too — same warehouse, different location.
  \textbf{The cloud is that warehouse. The internet is the road to it.}
\end{picturethis}

\sectionrule

% --- What can the cloud do? ---
\section*{What Can the Cloud Do?}

\subsection*{Cloud Storage}

Cloud storage means saving your files on remote servers.

\begin{itemize}
  \item Photos you take on a phone can be automatically backed up to the cloud.
  \item If you lose or break your phone, the photos are still safe.
  \item You can look at those photos on a laptop, tablet, or a friend's phone.
\end{itemize}

\subsection*{Cloud Programmes}

Some programmes live in the cloud and run on servers, not your device.

\begin{itemize}
  \item You write a document in a browser tab — the document is stored on a cloud server.
  \item Multiple people can work on the same document at the same time, from different countries.
\end{itemize}

\subsection*{Cloud Computing}

Really powerful computing tasks can be done on cloud servers.

\begin{itemize}
  \item Scientists analysing huge amounts of data can rent cloud computing power.
  \item Games can run on cloud servers and stream to your screen like a video.
\end{itemize}

\sectionrule

% --- The cloud is real buildings ---
\section*{The Cloud Is Very Much on the Ground}

\character{Cloud Pal}{I know — calling it a cloud is a bit misleading.
  Really, I am data centres — big buildings full of servers, cooling systems, and cables.
  The clouds in the diagrams are just a shorthand for saying
  ``somewhere on the internet — we don't need to know exactly where.'' }

Big companies run enormous cloud services around the world.
These data centres consume a lot of electricity and water for cooling,
so engineers are always working on ways to make them more efficient and environmentally friendly.

\sectionrule

% --- Benefits and things to know ---
\section*{Why the Cloud Is Useful}

\begin{quickmap}
  \textbf{Cloud benefits at a glance:}
  \begin{itemize}
    \item \textbf{Access anywhere} — your files are available on any device.
    \item \textbf{Backup and safety} — if your device breaks, your data is not lost.
    \item \textbf{Sharing} — easily share files or work on them together.
    \item \textbf{No need for huge storage} — your device does not need to hold everything.
  \end{itemize}
\end{quickmap}

\begin{diagram}[Devices connect through the internet to cloud servers and data centres.]
  \resizebox{\linewidth}{!}{%
  \begin{tikzpicture}[>=Stealth]
    \node[draw=WarmOrange!80!black, rounded corners, fill=WarmOrange!20, minimum width=1.9cm, minimum height=0.8cm] (phone) at (-4.2,2.2) {Phone};
    \node[draw=InternetBlue, rounded corners, fill=SkyBlue!20, minimum width=2.1cm, minimum height=0.8cm] (laptop) at (-1.8,2.2) {Laptop};
    \node[draw=LeafGreen!80!black, rounded corners, fill=LeafGreen!20, minimum width=1.9cm, minimum height=0.8cm] (tablet) at (0.6,2.2) {Tablet};
    \node[draw=SoftPurple!80!black, rounded corners, fill=SoftPurple!16, minimum width=1.9cm, minimum height=0.8cm] (watch) at (3.0,2.2) {Smartwatch};

    \node[cloud,draw=InternetBlue!80!black,fill=InternetBlue!12,minimum width=5.6cm,minimum height=2.7cm,cloud puffs=14,cloud puff arc=110] (cloud) at (-0.7,0.1) {};
    \node at (-1.1,0.85) {\small\bfseries Cloud};
    \node[draw=DarkText, rounded corners, fill=LightGray, minimum width=0.8cm, minimum height=0.42cm] at (-1.9,-0.05) {};
    \node[draw=DarkText, rounded corners, fill=LightGray, minimum width=0.8cm, minimum height=0.42cm] at (-1.0,-0.05) {};
    \node[draw=DarkText, rounded corners, fill=LightGray, minimum width=0.8cm, minimum height=0.42cm] at (-0.1,-0.05) {};
    \node[draw=DarkText, rounded corners, fill=LightGray, minimum width=0.8cm, minimum height=0.42cm] at (0.8,-0.05) {};

    \node[draw=CoralRed!80!black, rounded corners, fill=CoralRed!16, minimum width=2.4cm, minimum height=1cm] (dc) at (4.8,0.1) {data centre};

    \draw[thick,->,InternetBlue] (phone) -- (cloud);
    \draw[thick,->,InternetBlue] (laptop) -- (cloud);
    \draw[thick,->,InternetBlue] (tablet) -- (cloud);
    \draw[thick,->,InternetBlue] (watch) -- (cloud);
    \draw[thick,->,InternetBlue] (cloud) -- node[above]{via the Internet} (dc);
  \end{tikzpicture}%
  }
\end{diagram}

\subsection*{Things to Think About}

\begin{itemize}
  \item You need an internet connection to access cloud files.
  \item You are trusting a company to keep your files safe and private.
  \item It is good to know that the cloud is not magic — it is real physical computers
        that need electricity, maintenance, and protection.
\end{itemize}

\sectionrule

\begin{newwords}
  \begin{description}[labelwidth=3.5cm, leftmargin=3.7cm]
    \item[\term{The cloud}]       Servers accessed via the internet, used to store data or run programmes remotely.
    \item[\term{Cloud storage}]   Saving files on remote servers so they can be reached from any device.
    \item[\term{Backup}]          A copy of data stored safely somewhere else in case the original is lost.
    \item[\term{Streaming}]       Receiving and using data (like a video or game) directly from a server as you need it.
    \item[\term{Remote server}]   A server that is in a different physical location from the user.
  \end{description}
\end{newwords}

\sectionrule

\begin{tryit}
  \textbf{On the device or in the cloud?}

  For each of the following, mark whether it is likely stored \textbf{on your device}
  or \textbf{in the cloud}:

  \begin{center}
    \begin{tabularx}{\linewidth}{Ycc}
      \textbf{Thing} & \textbf{On device?} & \textbf{In the cloud?} \\
      \hline
      A photo automatically backed up & $\circ$ & $\circ$ \\
      An app you downloaded & $\circ$ & $\circ$ \\
      A shared document you edit in a browser & $\circ$ & $\circ$ \\
      A song saved for offline listening & $\circ$ & $\circ$ \\
      A video you are streaming right now & $\circ$ & $\circ$ \\
    \end{tabularx}
  \end{center}

  \emph{(Hint: many things can be both!  A song might be stored on your device AND in the cloud.)}
\end{tryit}

\sectionrule

\begin{bigidea}
  \textbf{The cloud} is not floating in the sky —
  it is servers in buildings on the ground, connected to you through the internet.
  It lets you store, share, and access your files and programmes from anywhere.
\end{bigidea}

\character{Cloud Pal}{You have learned so much on this journey!
  The last chapter is all about \textbf{you} — how to be a smart, safe, and curious explorer
  in this amazing internet world.  See you at the finish line!}
